% Bagian: sets-functions-relations
% Bab: functions
% Subbagian: inverses

\documentclass[../../../include/open-logic-section]{subfiles}

\begin{document}

\olfileid[id]{sfr}{fun}{inv}
\olsection{Invers Fungsi}

\begin{explain}
Kita memandang fungsi sebagai pemetaan. Karena itu, pertanyaan yang wajar
tentang fungsi ialah apakah pemetaannya dapat ``dibalik.'' Sebagai contoh,
fungsi penerus $f(x) = x + 1$ dapat dibalik, dalam arti bahwa fungsi $g(y) = y
- 1$ ``meniadakan'' apa yang dilakukan oleh $f$.

Namun, kita harus berhati-hati. Walaupun definisi~$g$ menentukan suatu fungsi
$\Int \to \Int$, definisi itu tidak menentukan sebuah \emph{fungsi} $\Nat \to
\Nat$, karena $g(0) \notin \Nat$. Jadi, bahkan dalam kasus sederhana, tidak
sepenuhnya jelas apakah suatu fungsi dapat dibalik; jawabannya dapat bergantung
pada domain dan kodomainnya.

Hal ini dirumuskan lebih tepat melalui gagasan invers suatu fungsi.
\end{explain}

\begin{defn}
Suatu fungsi $g \colon B \to A$ adalah \emph{invers} dari fungsi $f \colon A
\to B$ jika $f(g(y)) = y$ dan $g(f(x)) = x$ untuk semua $x \in A$ dan $y \in
B$.
\end{defn}

Jika $f$ mempunyai invers~$g$, kita sering menulis $f^{-1}$ sebagai ganti~$g$.

\begin{explain}
Sekarang kita akan menentukan kapan suatu fungsi mempunyai invers. Calon yang
baik untuk invers dari $f\colon A \to B$ adalah $g\colon B \to A$ yang
``didefinisikan oleh''
\[
g(y) = \text{``satu-satunya'' $x$ sedemikian sehingga $f(x) = y$.}
\]
Namun, tanda petik pada ``didefinisikan oleh'' (dan ``satu-satunya'')
menunjukkan bahwa ini belum tentu merupakan definisi. Setidaknya, cara ini
tidak selalu berhasil secara umum. Agar definisi tersebut menentukan sebuah
fungsi, harus ada satu dan hanya satu~$x$ sehingga $f(x) = y$---keluaran~$g$
harus ditentukan secara unik. Selain itu, keluaran tersebut harus ditentukan
untuk setiap $y \in B$. Jika terdapat $x_1$ dan $x_2 \in A$ dengan $x_1 \neq
x_2$, tetapi $f(x_1) = f(x_2)$, maka $g(y)$ tidak ditentukan secara unik untuk
$y = f(x_1) = f(x_2)$. Jika sama sekali tidak terdapat~$x$ sehingga $f(x) =
y$, maka $g(y)$ sama sekali tidak ditentukan. Dengan kata lain, agar $g$
terdefinisi, $f$ harus sekaligus !!{injective} dan !!{surjective}.

Mari kita telaah secara perlahan. Kita bagi pertanyaannya menjadi dua. Jika
diberikan fungsi~$f\colon A \to B$, kapan terdapat fungsi $g\colon B \to A$
sehingga $g(f(x)) = x$? Fungsi $g$ semacam itu ``meniadakan'' apa yang
dilakukan oleh $f$, dan disebut \emph{invers kiri} dari~$f$. Kedua, kapan
terdapat fungsi $h\colon B \to A$ sehingga $f(h(y)) = y$? Fungsi $h$ semacam
itu disebut \emph{invers kanan} dari~$f$---$f$ ``meniadakan'' apa yang
dilakukan oleh~$h$.
\end{explain}

\begin{prop}
Jika $f\colon A \to B$ bersifat !!{injective} dan domainnya tidak kosong, maka
terdapat \emph{invers kiri}~$g\colon B \to A$ dari~$f$ sehingga
$g(f(x)) = x$ untuk semua $x \in A$.
\end{prop}

\begin{proof}
Andaikan $f\colon A \to B$ bersifat !!{injective} dan domainnya tidak kosong.
Tinjau suatu $y \in B$.
Jika $y \in \ran{f}$, terdapat $x \in A$ sehingga $f(x) = y$. Karena $f$
bersifat !!{injective}, hanya ada satu~$x \in A$ semacam itu. Jadi, kita dapat
mendefinisikan $g(y) = x$, yaitu $g(y)$ adalah ``satu-satunya'' $x \in A$
sedemikian sehingga $f(x) = y$. Jika $y \notin \ran{f}$, karena domainnya
tidak kosong, kita dapat memetakannya ke sembarang~$a \in A$. Jadi, kita dapat
memilih $a \in A$ dan
mendefinisikan $g\colon B \to A$ dengan:
\[
g(y) = \begin{cases}
    x & \text{jika $f(x) = y$}\\
    a & \text{jika $y \notin \ran{f}$.}
\end{cases}
\]
Fungsi ini terdefinisi untuk semua $y \in B$, karena untuk setiap $y \in
\ran{f}$ terdapat tepat satu $x \in A$ sedemikian sehingga $f(x) = y$.
Berdasarkan definisi, jika $y = f(x)$, maka $g(y) = x$, yaitu $g(f(x)) = x$.
\end{proof}

\begin{prob}
Tunjukkan bahwa jika $f\colon A \to B$ mempunyai invers kiri~$g$, maka
$f$~bersifat !!{injective}.
\end{prob}

\begin{prop}
    Jika $f\colon A \to B$ bersifat !!{surjective}, ada
    \emph{invers kanan}~$h\colon B \to A$ dari~$f$ yang memenuhi $f(h(y)) = y$
    untuk semua~$y \in B$.
\end{prop}

\begin{proof}
Andaikan $f\colon A \to B$ bersifat !!{surjective}. Tinjau suatu $y \in B$.
Karena $f$~bersifat !!{surjective}, terdapat $x_y \in A$ dengan $f(x_y) = y$.
Lalu kita dapat mendefinisikan $h(y) = x_y$, yaitu untuk setiap $y \in B$ kita
memilih suatu $x \in A$ sehingga $f(x) = y$; karena $f$~bersifat
!!{surjective}, selalu ada sekurang-kurangnya satu yang dapat
dipilih.\footnote{Karena $f$ bersifat !!{surjective}, untuk setiap~$y \in B$
himpunan $\Setabs{x}{f(x) = y}$ tidak kosong. Definisi~$h$ mengharuskan kita
memilih satu $x$ dari setiap himpunan tersebut. Bahwa pilihan ini selalu dapat
dilakukan sebenarnya tidaklah jelas---kemungkinan untuk membuat semua pilihan
ini begitu saja diasumsikan sebagai suatu aksioma. Dengan kata lain,
proposisi ini mengasumsikan apa yang disebut Aksioma Pilihan, suatu persoalan
yang akan kita \oliflabeldef{sth:choice::chap}{tinjau kembali dalam
\olref[sth][choice][]{chap}}{abaikan untuk sementara}. Namun, dalam banyak kasus
khusus, misalnya ketika $A = \Nat$ atau berhingga, atau ketika $f$
bersifat !!{bijective}, Aksioma Pilihan tidak diperlukan. (Dalam kasus khusus
ketika $f$ bersifat !!{bijective}, untuk setiap $y \in B$ himpunan
$\Setabs{x}{f(x) = y}$ mempunyai tepat satu !!{element}, sehingga tidak ada
pilihan yang perlu dibuat.)}
Berdasarkan definisi, jika $x = h(y)$, maka $f(x) = y$, yaitu untuk setiap $y
\in B$, $f(h(y)) = y$.
\end{proof}

\begin{prob}
Tunjukkan bahwa jika $f\colon A \to B$ mempunyai invers kanan~$h$, maka
$f$~bersifat !!{surjective}.
\end{prob}

\begin{explain}
  Dengan menggabungkan gagasan dalam pembuktian sebelumnya, sekarang kita
  memperoleh bahwa setiap !!{bijection} mempunyai invers, yaitu terdapat satu
  fungsi yang sekaligus merupakan invers kiri dan invers kanan dari~$f$.
\end{explain}

\begin{prop}\ollabel{prop:bijection-inverse}
Jika $f\colon A \to B$ bersifat !!{bijective}, ada
fungsi~$f^{-1}\colon B \to A$. Untuk semua $x \in A$, fungsi ini memenuhi
$f^{-1}(f(x)) = x$; dan untuk semua $y \in B$, $f(f^{-1}(y)) = y$.
\end{prop}

\begin{proof}
Latihan.
\end{proof}

\begin{prob}
Buktikan \olref[sfr][fun][inv]{prop:bijection-inverse}. Anda harus
mendefinisikan~$f^{-1}$, menunjukkan bahwa ia merupakan suatu fungsi, dan
menunjukkan bahwa ia adalah invers dari~$f$, yaitu $f^{-1}(f(x)) = x$ dan
$f(f^{-1}(y)) = y$ untuk semua $x \in A$ dan $y \in B$.
\end{prob}

\begin{explain}
Ada cara yang sedikit lebih umum untuk memperoleh invers. Dalam
\olref[kin]{sec}, kita melihat bahwa setiap fungsi $f$ menginduksi
!!a{surjection} $f' \colon A \to \ran{f}$ dengan menetapkan $f'(x) = f(x)$
untuk semua $x \in A$. Jelas bahwa jika $f$~bersifat !!{injective}, maka
$f'$~bersifat !!{bijective}, sehingga mempunyai invers yang unik menurut
\olref{prop:bijection-inverse}. Dengan sedikit penyalahgunaan notasi, kita
kadang-kadang menyebut invers dari $f'$ semata-mata sebagai ``invers dari~$f$.''
\end{explain}

\begin{prop}\ollabel{prop:left-right}%
  Tunjukkan bahwa jika $f\colon A \to B$ mempunyai invers kiri~$g$ dan invers
  kanan~$h$, maka $h = g$.
\end{prop}

\begin{proof}
  Latihan.
\end{proof}

\begin{prob}
  Buktikan \olref[sfr][fun][inv]{prop:left-right}.
\end{prob}

\begin{prop}\ollabel{prop:inverse-unique}
Setiap fungsi~$f$ mempunyai paling banyak satu invers.
\end{prop}

\begin{proof}
  Andaikan $g$ dan $h$ keduanya merupakan invers dari~$f$. Maka, secara
  khusus, $g$~merupakan invers kiri dari~$f$ dan $h$~merupakan invers kanan.
  Menurut \olref{prop:left-right}, $g = h$.
\end{proof}

\end{document}
